\section{本讲小结}

求电场先确定源电荷分布与场点，利用方向、对称和叠加减少积分。场强是矢量，受力还需乘带符号的受力电荷。电场线表示局部方向与通量密度，不能替代粒子轨迹，也不能单凭疏密确定电势。

高斯定律普遍成立，直接求场还需使通量积分能够化为场强乘面积的对称条件。球壳、实心球和无限厚板分别给出球内零场、线性径向场和线性板内场；补形法则把缺口和空腔化为完整分布的叠加。

偶极远场应保留首个不抵消的小参数项。分布电荷受力必须排除自身场，平板吸引力和导体表面压力都因此含有因子 $1/2$。涉及平衡和运动时，还应说明机械约束、稳定方向与近似范围，不能把零合力和稳定混同。

